\section{Biểu đồ Use Case tổng thể}

Biểu đồ Use Case tổng thể được sử dụng để mô tả các tác nhân tham gia hệ thống và các chức năng chính mà từng tác nhân có thể thực hiện. Dựa trên các chức năng đã triển khai, hệ thống KMA Chatbot tuyển sinh có ba nhóm tác nhân chính: người dùng hỏi đáp, người quản trị nội dung và người quản trị hệ thống.

Trong đó, người dùng hỏi đáp tương tác trực tiếp với chatbot để đặt câu hỏi tuyển sinh và nhận phản hồi từ Rasa hoặc RAG. Người quản trị nội dung quản lý tài liệu, biểu mẫu, tài liệu ngữ cảnh RAG và dữ liệu phục vụ hỏi đáp. Người quản trị hệ thống chịu trách nhiệm quản lý tài khoản, vai trò, quyền truy cập và theo dõi các báo cáo thống kê của hệ thống.

\subsection{Danh sách tác nhân}
\begin{table}[H]
\centering
\caption{Danh sách tác nhân trong hệ thống}
\label{tab:actors_list}
\begin{tabular}{|p{3.5cm}|p{6.5cm}|p{5cm}|}
\hline
\textbf{Tác nhân} & \textbf{Mô tả} & \textbf{Phạm vi sử dụng} \\ \hline
Người dùng hỏi đáp & Là người sử dụng chatbot để tra cứu thông tin tuyển sinh, ngành học, học phí, thủ tục, văn bản và biểu mẫu & Hỏi đáp với Rasa, hỏi đáp với RAG \\ \hline
Người quản trị nội dung & Là người phụ trách cập nhật, quản lý và kiểm soát nguồn dữ liệu phục vụ chatbot & Quản lý tài liệu biểu mẫu, tài liệu ngữ cảnh RAG, dữ liệu huấn luyện và câu hỏi gợi ý \\ \hline
Người quản trị hệ thống & Là người quản lý vận hành, phân quyền và giám sát hệ thống & Quản lý người dùng, vai trò, quyền truy cập và thống kê hệ thống \\ \hline
\end{tabular}
\end{table}

\subsection{Danh sách Use Case tổng thể}

Dựa trên chức năng đã triển khai trong hệ thống, các Use Case chính được xác định như sau:

\begin{table}[H]
\centering
\caption{Danh sách Use Case tổng thể}
\label{tab:usecase_list}
\begin{tabular}{|p{1.5cm}|p{3.5cm}|p{3cm}|p{7cm}|}
\hline
\textbf{Mã Use Case} & \textbf{Tên Use Case} & \textbf{Tác nhân chính} & \textbf{Mô tả} \\ \hline
UC-01 & Hỏi đáp thường với Rasa & Người dùng hỏi đáp & Người dùng đặt câu hỏi phổ biến, câu hỏi theo luồng định sẵn hoặc nghiệp vụ có cấu trúc \\ \hline
UC-02 & Hỏi đáp RAG theo tài liệu & Người dùng hỏi đáp & Người dùng đặt câu hỏi và nhận câu trả lời dựa trên tài liệu đã upload, ingest và chia chunks \\ \hline
UC-03 & Chọn tài liệu ngữ cảnh khi hỏi RAG & Người dùng hỏi đáp & Người dùng chọn tài liệu cụ thể để làm ngữ cảnh cho quá trình hỏi đáp RAG \\ \hline
UC-04 & Quản lý tài liệu biểu mẫu & Người quản trị nội dung & Quản lý tài liệu tại /docs, gồm tìm kiếm, lọc, tạo tài liệu, theo dõi loại file và dung lượng \\ \hline
UC-05 & Quản lý tài liệu ngữ cảnh RAG & Người quản trị nội dung & Quản lý tài liệu tại /context-docs, gồm upload, scan, retry pipeline và theo dõi trạng thái xử lý \\ \hline
UC-06 & Theo dõi trạng thái xử lý tài liệu & Người quản trị nội dung & Theo dõi trạng thái processed, pending, failed và số chunks của tài liệu \\ \hline
UC-07 & Quản lý câu hỏi gợi ý và dữ liệu huấn luyện & Người quản trị nội dung & Bổ sung hoặc điều chỉnh dữ liệu phục vụ chatbot/Rasa \\ \hline
UC-08 & Xem báo cáo thống kê & Người quản trị hệ thống & Xem thống kê người dùng, hội thoại, chatbot, NLP và tài liệu tại nhóm /statistics/* \\ \hline
UC-09 & Quản lý vai trò & Người quản trị hệ thống & Quản lý vai trò tại /roles \\ \hline
UC-10 & Quản lý quyền & Người quản trị hệ thống & Quản lý quyền sử dụng chức năng tại /permissions \\ \hline
UC-11 & Quản lý người dùng & Người quản trị hệ thống & Quản lý danh sách tài khoản, vai trò, trạng thái và thao tác quản trị tại /users \\ \hline
\end{tabular}
\end{table}

\subsection{Mô tả quan hệ giữa các Use Case}

Trong hệ thống, một số Use Case có quan hệ mở rộng hoặc bao gồm nhau.

Use Case Hỏi đáp RAG theo tài liệu có liên quan đến Use Case Chọn tài liệu ngữ cảnh khi hỏi RAG, vì khi sử dụng RAG Chat, người dùng có thể lựa chọn tài liệu cụ thể để giới hạn phạm vi truy xuất. Điều này giúp câu trả lời tập trung hơn vào tài liệu mà người dùng quan tâm.

Use Case Quản lý tài liệu ngữ cảnh RAG bao gồm các thao tác như upload tài liệu, scan tài liệu mới, retry pipeline và theo dõi trạng thái xử lý. Vì vậy, Use Case Theo dõi trạng thái xử lý tài liệu có thể được xem là một phần quan trọng trong quản lý tài liệu ngữ cảnh.

Use Case Xem báo cáo thống kê phục vụ người quản trị hệ thống trong việc giám sát hoạt động. Các thống kê có thể bao gồm thống kê người dùng, hội thoại, chatbot, NLP và tài liệu.

Use Case Quản lý vai trò, Quản lý quyền và Quản lý người dùng thuộc nhóm quản trị hệ thống. Các Use Case này có mối quan hệ chặt chẽ vì quyền truy cập của người dùng phụ thuộc vào vai trò được gán và các quyền tương ứng.

\subsection{Biểu đồ Use Case tổng thể}

\begin{figure}[H]
    \centering
    \includegraphics[width=\textwidth]{source/usecase_tong_the.png}
    \caption{Biểu đồ Use Case tổng thể của hệ thống KMA Chatbot tuyển sinh}
    \label{fig:overall_use_case}
\end{figure}
